\documentclass[12pt]{article}
\usepackage[utf8]{inputenc}
\usepackage[T1]{fontenc}
\usepackage[french]{babel}
\usepackage{amsmath}
\usepackage{chemfig}
\usepackage{mol2chemfig}
\usepackage{float}
\usepackage{hyperref}

\title{Résumé : Nomenclature des Molécules Organiques}
\author{}
\date{}

\begin{document}

\maketitle

\section{Nomenclature des Alcanes}

\subsection{Alcanes Linéaires}
Les alcanes sont des hydrocarbures saturés, c'est-à-dire qu'ils ne contiennent que des liaisons simples entre atomes de carbone et d'hydrogène. Leur nom dépend du nombre d'atomes de carbone dans la chaîne principale.

\begin{center}
\begin{tabular}{|c|c|c|}
\hline
Nombre de Carbone & Formule brute & Nom de l'alcane linéaire \\ \hline
1 & $\mathrm{CH_4}$ & Méthane \\ \hline
2 & $\mathrm{C_2H_6}$ & Éthane \\ \hline
3 & $\mathrm{C_3H_8}$ & Propane \\ \hline
4 & $\mathrm{C_4H_{10}}$ & Butane \\ \hline
5 & $\mathrm{C_5H_{12}}$ & Pentane \\ \hline
6 & $\mathrm{C_6H_{14}}$ & Hexane \\ \hline
7 & $\mathrm{C_7H_{16}}$ & Heptane \\ \hline
8 & $\mathrm{C_8H_{18}}$ & Octane \\ \hline
9 & $\mathrm{C_9H_{20}}$ & Nonane \\ \hline
10 & $\mathrm{C_{10}H_{22}}$ & Décane \\ \hline
\end{tabular}
\end{center}

\subsection{Numérotation des Atomes de Carbone}
Pour nommer correctement les alcanes ramifiés, il est nécessaire de numéroter les atomes de carbone de la chaîne principale. Voici les règles de base~:
\begin{itemize}
    \item Identifier la chaîne carbonée la plus longue.
    \item Numéroter les atomes de carbone de cette chaîne de manière à ce que les substituants aient les numéros les plus bas possibles.
\end{itemize}

\subsection{Numérotation des Atomes de Carbone}
Pour nommer correctement les alcanes ramifiés, il est nécessaire de numéroter les atomes de carbone de la chaîne principale. Voici les règles de base~:
\begin{itemize}
    \item Identifier la chaîne carbonée la plus longue.
    \item Numéroter les atomes de carbone de cette chaîne de manière à ce que les substituants aient les numéros les plus bas possibles.
\end{itemize}

Exemple: 3-R-heptane (R est ici un Substituant quelconque)
\begin{center}
    \chemfig{C_1-C_2-C_3(-[:90]R)-C_4-C_5-C_6-C_7}
\end{center}

\subsection{Alcanes Ramifiés}
Les alcanes ramifiés ont des substituants (groupes alkyle) attachés à la chaîne principale. Attention à bien choisir la chaîne la plus longue!

Les substituants sont nommés comme les alcanes, mais en remplaçant \textbf{-ane} par \textbf{-yle}.

\begin{center}
\begin{tabular}{|c|c|c|}
\hline
Substituant & Formule & Nom \\ \hline
Méthyle & $\mathrm{CH_3-}$ & Méthyl- \\ \hline
Éthyle & $\mathrm{C_2H_5-}$ & Éthyl- \\ \hline
Propyle & $\mathrm{C_3H_7-}$ & Propyl- \\ \hline
\end{tabular}
\end{center}

\subsubsection{Exemples de Nomenclature}
\begin{itemize}
    \item \textbf{2-Méthylbutane}~:
    \begin{center}
    \chemfig{
H_3C-[:330,,2]{CH}(-[:270,,,1]CH_3)%
-[:30]{CH_2}-[:330,,,1]CH_3}
    
    \end{center}

    \item \textbf{2,3-Diméthylpentane}~:
    \begin{center}
    \chemfig{@C(-[:90]CH_3)-[:30]C(-[:330]CH_3)-[:30]C-[:30]CH_3-[:30]CH_3}
    \end{center}
\end{itemize}

\subsection{Groupes Caractéristiques et Fonctions}
Les groupes caractéristiques déterminent la famille de la molécule organique. Voici quelques exemples de groupes caractéristiques et leurs fonctions associées~:

\begin{center}
\begin{tabular}{|c|c|c|}
\hline
Groupe Caractéristique & Fonction & Formule Générale \\ \hline
Groupe hydroxyle —OH & Alcool & R-OH \\ \hline
Groupe carbonyle —C— & Aldéhyde & R-C-H \\ \hline
 & Cétone & R1-C-R2 \\ \hline
Groupe carboxyle —C-OH & Acide Carboxylique & R-C-OH \\ \hline
Groupe amino —NH2 & Amine & R-NH2 \\ \hline
\end{tabular}
\end{center}

\subsection{Exemples de Noms de Molécules avec Fonctions}
\begin{itemize}
    \item \textbf{Propan-2-ol}~:
    \begin{center}
    \chemfig{@C(-[:90]OH)-[:30]CH_3-[:330]CH_3}
    \end{center}

    \item \textbf{Butan-2-one}~:
    \begin{center}
    \chemfig{@C(=[:90]O)-[:30]CH_3-[:330]CH_3-[:30]CH_3}
    \end{center}
\end{itemize}

\subsection{Exercices}
\begin{enumerate}
    \item Nommez les molécules suivantes~:
    \begin{itemize}
        \item \chemfig{@C(-[:90]CH_3)-[:30]CH_2-[:30]CH_2-[:30]CH_3}
        \item \chemfig{@C(-[:90]CH_3)-[:30]C(-[:330]CH_3)-[:30]CH_3}
    \end{itemize}

    \item Dessinez les formules développées des molécules suivantes~:
    \begin{itemize}
        \item 2,2-Diméthylbutane
        \item 3-Éthylhexane
    \end{itemize}
\end{enumerate}


\end{document}
